\chapter*{Preface}
\addcontentsline{toc}{chapter}{Preface}

Many people first meet \LaTeX{} at an awkward moment. A thesis has grown too
large for comfortable manual formatting. An equation has become difficult to
edit. A journal has supplied a template full of unfamiliar commands. The
learner is expected to produce a polished document while also trying to
understand a new way of writing.

This book begins earlier and moves more gently. It assumes that you may never
have seen a source file, a command, or a compiler. In the first lessons we use
the smallest document that \LaTeX{} can compile. We then add one useful idea at a
time, inspect the result, and return to improve the same scholarly project.
The aim is not to memorise hundreds of commands. It is to understand enough
structure to make sound decisions, find trustworthy documentation, and repair
problems calmly.

Our continuing project is an article called ``A Reproducible Comparison of
Two Measurement Methods''. Its values are simulated demonstration data, not
real findings. That neutral example lets us practise the architecture of a
scientific article without pretending to conduct a study. Shorter examples
from biology, chemistry, physics, mathematics, earth science, computing,
engineering, health, and interdisciplinary work show how the same principles
travel across fields.

You will see errors in these pages. Some are deliberate. A missing brace or a
misspelt command is not evidence that you are unsuited to \LaTeX{}; it is a clue
about what the compiler could not understand. We shall learn to read the first
useful clue, examine the nearby source, and test a small correction.

Three professional outcomes guide the journey: a scholarly manuscript, a
thesis or dissertation, and a technical or scientific report. Each requires
more than attractive pages. It must also have maintainable source, reliable
references, accessible figures and tables, clear metadata, and a reproducible
build. By the end, you should be able to create those documents, explain their
structure, and prepare clean files for another person to compile.

Software changes, and publisher or institutional instructions differ. This
book therefore distinguishes lasting \LaTeX{} ideas from interface details and
local rules. When a current instruction controls a submission, follow that
instruction. When you are free to choose, prefer a simple, semantic, and
portable solution.
